\chapter{Lo que estas decisiones tienen en común}
\label{cap:d-patrones}
Leídas juntas, las decisiones de las 12 fases siguen cinco reglas. Son las que conviene defender en el coloquio, porque
explican \emph{cómo} se trabajó y no solo \emph{qué} salió.

\section{1. Ninguna cifra se inventó}
Cada vez que faltó un dato, se eligió entre declararlo hueco o probarlo como supuesto con varios valores. Nunca se rellenó.
\begin{center}
\small
\begin{tabularx}{\linewidth}{>{\raggedright\arraybackslash}p{0.35\linewidth}Y}
\encabezado \h{Dato que faltaba} & \h{Cómo se trató} \\
Ocupación del sur desde 2025 & Hueco; el sur se mide con presión de llegada (3.4) \\
\filagris Conversión de Facebook para turismo & Supuesto con barrido: 3, 5.75 y 6.38\,\% (6.3) \\
Golpe de una tormenta & Supuesto con barrido: 0, 25 y 50\,\% (5.6) \\
\filagris Efecto de la campaña en los escenarios & No se metió un número sin fuente (5.7) \\
Feriados sin tipo de cambio & Vacíos, sin copiar el día anterior (2.7) \\
\filagris Reseñas y fotos de platillos del sur & No se muestran; las aportará el equipo (P.4) \\
\bottomrule
\end{tabularx}
\normalsize
\end{center}

\section{2. Cada modelo se comparó contra una regla simple}
El Radar contra ``igual que este mes'' (130 contra 126); el pronóstico contra ``el mismo mes del año pasado'' (en Cancún ganó
la regla simple, y se usa); Markov contra la persistencia (mejor puntaje a 1, 4 y 8 semanas). Siempre se dice cuánto le gana
cada modelo, aunque sea poco.

\section{3. Los criterios se fijaron antes de ver el ganador}
Brandon eligió \emph{cómo} elegir (``el que más acierta y más cambios anticipa'', ``menor error con rango $\geq80\,\%$'') y el
código aplica el criterio solo. Cuando la auditoría cambió los datos, el modelo ganador del Radar cambió de Random Forest a
regresión logística, y se respetó el criterio, no el nombre.

\section{4. Cuando algo falló, se contó}
\begin{itemize}
\item El índice que daba a Cancún como tranquilo en julio (4.5).
\item La ``ocupación'' de 2,560,068 y los visitantes contados dos veces (Fase 4).
\item Los conteos inflados del DENUE y del Censo (Fase 1).
\item El corte de clima raro que marcaba el 47\,\% de las semanas (7.5).
\item Todo el dinero de la Ruta en junio (6.5).
\item Junio como mes recomendado, siendo el más lluvioso (5.11).
\item Las fotos de Mérida y Campeche (P.3).
\end{itemize}
Cada uno se corrigió con datos y quedó escrito. Varios se convirtieron en pruebas automáticas para que no vuelvan a pasar.

\section{5. La redistribución tiene reglas que no se pueden romper}
La campaña no anuncia en temporada alta, no rebasa la capacidad probada, reparte entre los tres lugares, no depende de una
sola plataforma, se pausa con mal clima y manda al sur a quien iba al norte lleno. Y se midió el precio de cuidar: \textbf{cero
visitantes}.
